\section{Background and Related Work}
\label{sec:background}

\subsection{MSI data and preprocessing}
An MSI experiment rasters a tissue section and records a spectrum at each
location, giving a data cube with two spatial dimensions and one $m/z$
dimension~\citep{caprioli1997maldi,buchberger2018mass}. Spectra vary in total
intensity between pixels for technical reasons, so they are usually
normalised, most often by total ion current (TIC), before
comparison~\citep{deininger2011normalization}. Spectra may be kept as fine
$m/z$ bins or reduced to centroided peaks. Because neighbouring bins of one
molecular peak are strongly correlated, any evaluation that counts bins can
reward several bins of the same peak. This motivates the peak-level scoring
used in this report (\cref{sec:evaluation}). Tools such as
Cardinal~\citep{bemis2015cardinal} and M\textsuperscript{2}aia~\citep{cordes2021m2aia}
provide standard data handling.
For our work, the key consequence is that every method must receive the same
input and be scored at the same resolution, or differences in preprocessing
masquerade as differences in method.

\subsection{Peak selection}
Classical pipelines such as MALDIquant detect local maxima above a noise
threshold in each spectrum~\citep{gibb2012maldiquant}. Machine-learning peak
pickers learn what a peak looks like from annotated
spectra~\citep{li2022machinelearningpeakpicking}. Both operate spectrum by
spectrum and do not use the spatial arrangement of the tissue. Learned
representation methods instead rank $m/z$ bins by how a model trained on the
whole section uses them. This report compares such learned rankings with each
other and with simple model-free rankings, under one shared evaluation.

\subsection{Learned peak selection: msiPL and S3PL}
msiPL trains a fully connected VAE~\citep{kingma2014autoencoding} on
individual spectra, clusters the low-dimensional latent space, and reads peaks
from the network weights~\citep{abdelmoula2021peak}. Its companion work,
massNet, provides the glioblastoma (GBM) dataset used
here~\citep{abdelmoula2022massnet}. S3PL learns a 3D convolution over a
spectral patch around each pixel, applies an attention step to the centre
spectrum, reconstructs the full normalised patch, and selects the bins that
receive the most attention across pixels~\citep{weigand2026spatial}. S3PL also
introduced the mask-based spatial-correlation F1 score (mSCF1) used to judge
peak lists against expert annotations, and released the peak budgets and
protocols for GBM, colorectal adenocarcinoma (CAC) and renal cell carcinoma
(RCC). It reported 0.496, 0.601 and 0.470 mSCF1 on these collections, against
0.403, 0.431 and 0.159 for msiPL. Related representation-learning work in MSI
includes deep models of 3D data~\citep{inglese2017deep}, low-dimensional
ion-image embeddings~\citep{guo2024deepion} and IonMorphNet~\citep{weigand2026ionmorphnet}.
S3PL's result motivates our spatial question. Its convolution, objective,
normalisation and selector all differ from msiPL, so it is not clear which
ingredient produces the gain, or whether a much smaller amount of spatial
context would suffice.

\subsection{Representation quality versus downstream use}
A representation that reconstructs the data well is not necessarily the one
that serves a later task best. Mazarura, Jarvis and Bau compared autoencoders,
PCA and NMF on MSI data and found that reconstruction error, cluster quality
and downstream classification ranked the models
differently~\citep{mazarura2025autoencoder}. We test the same distinction for
peak ranking (\cref{sec:res-encoder}).

\subsection{Attribution and segment-then-correlate}
Integrated Gradients attributes a model output to input features by
integrating gradients along a path from a baseline input to the actual
input~\citep{sundararajan2017axiomatic}. Applied to a GMM posterior computed
from a VAE latent space, it ranks $m/z$ bins by how strongly they move a pixel
towards a tissue segment, through the full nonlinear encoder. An older and
simpler route is to segment the tissue first and then correlate each ion image
with the segments~\citep{alexandrov2010spatial,alexandrov2011efficient}. Our
pipeline contains both routes on the same segmentation, which allows a direct
comparison.

\subsection{Spatial statistics and spatially aware methods}
Ion images of biologically meaningful molecules tend to be spatially
structured rather than noisy. Moran's~$I$ measures spatial
autocorrelation~\citep{moran1950notes}, and spatial autocorrelation and
spatial-chaos measures have been used to judge the informativeness of MSI ion
images~\citep{alexandrov2013testing,cassese2016spatial}. Spatially aware
segmentation smooths features using neighbourhood weights before
clustering~\citep{alexandrov2011efficient}. Graph-based models of spatial
omics and MSI propagate information between neighbouring locations. These are
discussed as future work in \cref{sec:conclusion}. Our study places spatial
information at three points (the encoder input, the segmentation and the
ranking statistic) and tests each one under the same evaluation.

\subsection{Gap addressed}
Published learned peak-selection methods are usually evaluated against one or
two baselines, often with different peak counts, single training runs and no
model-free references. It has not been tested under a common budget whether
their advantages come from spatial input, from the learned representation, or
from the ranking step, nor how stable those advantages are across training
seeds. This report addresses that gap for a VAE-based pipeline, with
reproduced msiPL and S3PL baselines.
